\documentclass{zhvt-classic}

% 使用本地源流明体（GenRyuMin2 TW）按文件路径加载，不依赖系统字体名
% 版框（内框／界栏／书口侧栏／鱼尾）改用与版心印章同一印泥棕色：
%   印章图片 yshs.png 实色 = RGB(135,68,52) = HTML 874434
\usepackage{xcolor}
\usepackage{tikz}
% sealbrown 双注册：xcolor 供 TikZ/pgf，L3 供 cls 版框
\definecolor{sealbrown}{HTML}{874434}
\ExplSyntaxOn
\color_set:nnn { sealbrown } { HTML } { 874434 }
\ExplSyntaxOff

\zhvtset{
  grid color       = sealbrown,
  main font        = GenRyuMin2TW-SB,
  main font path   = C:/Users/Administrator/texmf/fonts/opentype/public/GenRyuMin2TW/,
jiazhu font      = GenRyuMin2TW-M,
  jiazhu font path = C:/Users/Administrator/texmf/fonts/opentype/public/GenRyuMin2TW/,
  % 页码框容纳 4 字：正常页「九十七」+前置空 1 格=4 行恰好；
  % 100+ 页「一百零一」共 5 行超出，末字「一/二/三」被框裁掉（2026-10-04 实测），故放宽到 5
  shukou page number chars = 5,
  % 书口书名框高从默认 6 字减到 4 字，把省出的 2 字高拨给 chapter 框，
  % 使 6 字书名+卷名放进 9 字高的 chapter 框后不再压线（2026-10-05）
  shukou title chars = 4,
}

\title{史记}[陈丞相世家]
\author{汉·司马迁}
\maker{闲叟漫记·重製}


% ===== 索隐述赞样式（2026-11-08 四次定稿 · 愚大人亲手改定参数）=====
% ⚠ 竖排前提：自定义字体家族必须带 RawFeature={vertical:+vert}；
%   simhei.ttf 加 +vhal 会整体加载失败静默回退，只写 +vert；simfang 同理。
\newCJKfontfamily{\shuzanfont}[Path=C:/Windows/Fonts/,RawFeature={vertical:+vert}]{simhei.ttf}
\newCJKfontfamily{\shuzanbody}[Path=C:/Windows/Fonts/,RawFeature={vertical:+vert}]{simfang.ttf}
% 题头：152x30pt 圆角 6pt 棕底白字；baseline (0,2) 使整框落入界行格子不压版框线；
% 白字节点 x 不补偿、y +8.2pt（像素级墨迹实测 dx/dy≈0，L/R 留白约 20pt）。
\newcommand{\shuzanhead}[1]{%
  \begin{tikzpicture}[baseline={(0,2)}]
    \fill[sealbrown, rounded corners=6pt]
      (-76pt,-15pt) rectangle (76pt,15pt);
    \node[text=white, inner sep=0pt, anchor=center] at (0pt,8.2pt)
      {\shuzanfont\fontsize{23}{23}\selectfont #1};
  \end{tikzpicture}%
}
% ★ 修正书口卷名位置：chapter 框起始位置改回 cls 默认值，避免全局影响。
%   只在每个卷首页通过 \juan 命令里的空命令微调该页位置（安全，不会跨页污染）。
\ExplSyntaxOn
\AtBeginDocument{
  \dim_set:Nn \l__zhvt_shukou_chapter_title_start_dim
    { \l__zhvt_shukou_title_height_dim + \l__zhvt_font_size_dim * 2 }
  \dim_set:Nn \l__zhvt_shukou_page_number_start_dim
    { \l__zhvt_shukou_chapter_title_start_dim + \l__zhvt_shukou_chapter_title_height_dim }
}
\ExplSyntaxOff

\begin{document}

% ★★ 版心（书口）排布覆写：把书名放到鱼尾之下（2026-10-04）
%
%   cls 默认版心自上而下：书名（title 静态框，在鱼尾之上）→ 鱼尾 → 篇名（chapter 动态框）→ 页码
%   本模板改为：          鱼尾（位置不动）→ 书名 → 篇名 → 页码分割线 → 页码
%
%   手法：把 cls 放在 title 静态框里的书名，改放进 chapter 动态框（鱼尾之下）；
%         title 静态框留空。书名取 \l__zhvt_title_str，即 \title{...} 的第一参数，
%         所以换书时只改 \title，不必在此重复写书名。
%
%   ⚠ 三条必须遵守的写法（2026-10-04 实测）：
%     ① 汉字必须用 \hbox{...} 包住。flowfram 会把框内容装进一个宽 14.4mm（= 一个 line_sep
%        = 书口栏宽）的纵向 minipage 并按该宽度折行，而整页靠 PDF /Rotate 90 显示、
%        框内文字本身不变向 —— 不包 \hbox 时「围炉夜话」4 字会被逐字折成 4 行，
%        在版面上就成了横排 4 字并横向溢出版框。
%     ② 竖排下**一个 \hbox 就是一整列**。书名与篇名要同处一列，必须放进同一个 \hbox；
%        分成两个 \hbox 时，后一个会被当成新的一列、甩到隔壁栏去
%        （实测「序」跑到书口栏外侧 14.4mm 的正文列上）。中间的 \quad 在框内是横向空白，
%        显示出来正好是书名与篇名之间的纵向间隔。
%     ③ ★ 内容盒要写成 **与框等宽的定宽盒** `\hbox to \l__zhvt_shukou_chapter_title_height_dim{…\hfil}`。
%        实测：chapter 框是把内容盒沿纵向**居中**摆放的，内容宽度一变，书名就跟着漂：
%        自然宽度的 \hbox（书名 38.4mm）落在 x=157.2mm，加篇名后（57.6mm）落在 x=152.4mm，
%        目录页（无篇名）与正文页的书名就差半格。做成等宽盒后，无论有没有篇名，
%        内容固定从框顶起排，实测两页均落在 x=142.8mm，不再漂。
\ExplSyntaxOn
% 把书口 chapter 框（书名+篇名）起始位置下移 1 字高，避免 4 字卷名顶到鱼尾
% ⚠ cls 的 \__zhvt_calculate_dimensions: 在 \AtBeginDocument 里执行，会重算该 dim；
%   必须把覆写挂到 begin document 之后，否则会被二次覆盖（2026-10-05 实测）。
% (旧方案已移到导言区，直接修改 chapter 框起始位置)
% 卷标题命令 #1=书口篇名(鱼尾) #2=正文卷名（含卷序，卷名内空格用 ~）
%
% ★★ 目录页两项修订（2026-10-05，愚大人指出）★★
%   ① 页码起点统一对齐：
%      原状 —— 目录项写成 \par\hspace{4em}卷名\hspace{4em}页码，页码起点 = 卷名实宽 + 4em。
%      卷名实宽随字数变：「卷一~梁惠王上」6 字 +~、「卷七~离娄上」5 字 +~，
%      差 1 个字位 → 卷七~卷十 4 条的页码整体比其余 10 条高 27.21pt（一个字），看着上下不齐。
%      改法 —— 卷名补全角空格（宽 = 1 字 = 27.215pt）补足到 **7 字符位**再写进目录，
%      14 条页码起点遂全部落在同一条竖线上（起点对齐；末字随页码字数自然参差）。
%      ⚠ 只对目录里的副本 \l_tmpc_tl 补，正文卷名仍用 \l_tmpa_tl 原样 —— 正文像素不受影响。
%   ② 卷名整体上移半字位：前导 \hspace 由 4em 改 3.5em。
%      实测原状目录块在版框内偏下（天头留白 139.6pt / 地脚留白 110.7pt），
%      上移 0.5em = 13.6pt 后上下基本均分，视觉居中。
\NewDocumentCommand { \juan } { mm } {
	\clearpage
	\tl_set:Nn \l__zhvt_chapter_title_tl { #1 }
	\tl_set:Nn \l_tmpa_tl { #2 }
	\tl_set_eq:NN \l_tmpc_tl \l_tmpa_tl
	\int_step_inline:nn { 7 - \str_count:n { #2 } }
	  { \tl_put_right:Nn \l_tmpc_tl { 　 } }
	\tl_set:Nx \l_tmpb_tl { \zhnum { page } }
	\tl_put_left:Nn \l_tmpb_tl { \hspace { 4em } }
	\tl_put_left:Nn \l_tmpb_tl { \l_tmpc_tl }
	\tl_put_left:Nn \l_tmpb_tl { \par \hspace { 3.5em } }
	\addtocontents { toc } { \l_tmpb_tl }
	\skip_horizontal:n { \l__zhvt_font_size_dim * 4 }
	\l_tmpa_tl
	\par
	\thispagestyle { empty }
}
\cs_set_protected:Npn \__zhvt_set_shukou_content: {
	% ★ 鱼尾上方留白（title 静态框）嵌入闲叟钤印：
	%   整页 /Rotate 90，故此处把图用 \rotatebox{90} 转正为竖章；
	%   尺寸卡在版心条幅内（书内约 7mm 宽 × 33.8mm 高，留余量不压鱼尾）。
	\setstaticcontents * { title } {
		\vbox_to_ht:nn { \l__zhvt_baseline_skip_dim } {
			\vskip 14.2pt
			\hbox_to_wd:nn { \l__zhvt_shukou_title_height_dim } {
				\hfil
				\rotatebox [ origin = c ] { 90 } { \includegraphics [ width = 7mm , height = 33.8mm ] { yshs.png } }
				\hfil
			}
		}
	}
% ★ 版心（书口）统一改法（2026-11-16 愚大人定）：18pt＋框内居中＋书口文字同版框色
%   ① 原 27.2pt 时 7 字长篇名（如「廉颇蔺相如列传」）末字溢出 chapter 框、压住页码；
%      改 \fontsize{18pt}{18pt}，前后 \hfil 令字串在框内居中（不随篇名字数漂移）。
%   ② \vbox{\vskip 8.2pt} 定量下移，抵消缩字号后栏宽方向 8.2pt 上偏（像素实测 y 中心 543.3 = 原版）。
%      ⚠ \raisebox / \vspace* 在 flowfram 框内无效；换字号须同步换补偿：16pt→10.1、19pt→7.4。
%   ③ chapter 与 page~number 均置 \color{sealbrown}（HTML 874434），与版框／界行同色；正文仍黑。
%   ⚠ 页码字位、钤印、鱼尾、界行、正文一律未动；五部合订本同此口径。
	\setdynamiccontents * { chapter }
	{ \color{sealbrown} \vbox { \vskip 8.2pt \hbox to \l__zhvt_shukou_chapter_title_height_dim
		{ \hfil \fontsize{18pt}{18pt}\selectfont \l__zhvt_title_str \quad \l__zhvt_chapter_title_tl \hfil } } }
	\setdynamiccontents * { page~number } { \color{sealbrown} \quad \zhnum { page } }
}
\__zhvt_set_shukou_content:
\ExplSyntaxOff

\maketitle

% 单篇，无多卷目录，不渲染目录页

\juan{陈丞相世家}{卷五十六~陈丞相世家第二十六}
\hspace*{2em}陈丞相平者。阳武户牖乡人也。少时家贫。好读书。有田三十亩。独与兄伯居。伯常耕田。纵平使游学。平为人长美色。人或谓陈平曰。﹁贫何食而肥若是。﹂其嫂嫉平之不视家生产。曰。﹁亦食糠覈耳。有叔如此。不如无有。﹂伯闻之。逐其妇而弃之。\\
\hspace*{2em}及平长。可娶妻。富人莫肯与者。贫者平亦耻之。久之。户牖富人有张负。张负女孙五嫁而夫辄死。人莫敢娶。平欲得之。邑中有丧。平贫。侍丧。以先往後罢为助。张负既见之丧所。独视伟平。平亦以故後去。负随平至其家。家乃负郭穷巷。以弊席为门。然门外多有长者车辙。张负归。谓其子仲曰。﹁吾欲以女孙予陈平。﹂张仲曰。﹁平贫不事事。一县中尽笑其所为。独柰何予女乎。﹂负曰。﹁人固有好美如陈平而长贫贱者乎。﹂卒与女。为平贫。乃假贷币以聘。予酒肉之资以内妇。负诫其孙曰。﹁毋以贫故。事人不谨。事兄伯如事父。事嫂如母。﹂平既娶张氏女。赍用益饶。游道日广。\\
\hspace*{2em}里中社。平为宰。分肉食甚均。父老曰。﹁善。陈孺子之为宰。﹂平曰。﹁嗟乎。使平得宰天下。亦如是肉矣。﹂\\
\hspace*{2em}陈涉起而王陈。使周市略定魏地。立魏咎为魏王。与秦军相攻於临济。陈平固已前谢其兄伯。从少年往事魏王咎於临济。魏王以为太仆。说魏王不听。人或谗之。陈平亡去。\\
\hspace*{2em}久之。项羽略地至河上。陈平往归之。从入破秦。赐平爵卿。项羽之东王彭城也。汉王还定三秦而东。殷王反楚。项羽乃以平为信武君。将魏王咎客在楚者以往。击降殷王而还。项王使项悍拜平为都尉。赐金二十溢。居无何。汉王攻下殷。项王怒。将诛定殷者将吏。陈平惧诛。乃封其金与印。使使归项王。而平身间行杖剑亡。渡河。船人见其美丈夫独行。疑其亡将。要中当有金玉宝器。目之。欲杀平。平恐。乃解衣裸而佐刺船。船人知其无有。乃止。\\
\hspace*{2em}平遂至修武降汉。因魏无知求见汉王。汉王召入。是时万石君奋为汉王中涓。受平谒。入见平。平等七人俱进。赐食。王曰。﹁罢。就舍矣。﹂平曰。﹁臣为事来。所言不可以过今日。﹂於是汉王与语而说之。问曰。﹁子之居楚何官。﹂曰。﹁为都尉。﹂是日乃拜平为都尉。使为参乘。典护军。诸将尽讙。曰。﹁大王一日得楚之亡卒。未知其高下。而即与同载。反使监护军长者。﹂汉王闻之。愈益幸平。遂与东伐项王。至彭城。为楚所败。引而还。收散兵至荥阳。以平为亚将。属於韩王信。军广武。\\
\hspace*{2em}绛侯。灌婴等咸谗陈平曰。﹁平虽美丈夫。如冠玉耳。其中未必有也。臣闻平居家时。盗其嫂。事魏不容。亡归楚。归楚不中。又亡归汉。今日大王尊官之。令护军。臣闻平受诸将金。金多者得善处。金少者得恶处。平。反覆乱臣也。原王察之。﹂汉王疑之。召让魏无知。无知曰。﹁臣所言者。能也。陛下所问者。行也。今有尾生。孝己之行而无益处於胜负之数。陛下何暇用之乎。楚汉相距。臣进奇谋之士。顾其计诚足以利国家不耳。且盗嫂受金又何足疑乎。﹂汉王召让平曰。﹁先生事魏不中。遂事楚而去。今又从吾游。信者固多心乎。﹂平曰。﹁臣事魏王。魏王不能用臣说。故去事项王。项王不能信人。其所任爱。非诸项即妻之昆弟。虽有奇士不能用。平乃去楚。闻汉王之能用人。故归大王。臣裸身来。不受金无以为资。诚臣计画有可采者。大王用之。使无可用者。金具在。请封输官。得请骸骨。﹂汉王乃谢。厚赐。拜为护军中尉。尽护诸将。诸将乃不敢复言。\\
\hspace*{2em}其後。楚急攻。绝汉甬道。围汉王於荥阳城。久之。汉王患之。请割荥阳以西以和。项王不听。汉王谓陈平曰。﹁天下纷纷。何时定乎。﹂陈平曰。﹁项王为人。恭敬爱人。士之廉节好礼者多归之。至於行功爵邑。重之。士亦以此不附。今大王慢而少礼。士廉节者不来。然大王能饶人以爵邑。士之顽钝嗜利无耻者亦多归汉。诚各去其两短。袭其两长。天下指麾则定矣。然大王恣侮人。不能得廉节之士。顾楚有可乱者。彼项王骨鲠之臣亚父。锺离眜。龙且。周殷之属。不过数人耳。大王诚能出捐数万斤金。行反间。间其君臣。以疑其心。项王为人意忌信谗。必内相诛。汉因举兵而攻之。破楚必矣。﹂汉王以为然。乃出黄金四万斤。与陈平。恣所为。不问其出入。\\
\hspace*{2em}陈平既多以金纵反间於楚军。宣言诸将锺离眛等为项王将。功多矣。然而终不得裂地而王。欲与汉为一。以灭项氏而分王其地。项羽果意不信锺离眜等。项王既疑之。使使至汉。汉王为太牢具。举进。见楚使。即详惊曰。﹁吾以为亚父使。乃项王使。﹂复持去。更以恶草具进楚使。楚使归。具以报项王。项王果大疑亚父。亚父欲急攻下荥阳城。项王不信。不肯听。亚父闻项王疑之。乃怒曰。﹁天下事大定矣。君王自为之。原请骸骨归。﹂归未至彭城。疽发背而死。陈平乃夜出女子二千人荥阳城东门。楚因击之。陈平乃与汉王从城西门夜出去。遂入关。收散兵复东。\\
\hspace*{2em}其明年。淮阴侯破齐。自立为齐王。使使言之汉王。汉王大怒而骂。陈平蹑汉王。汉王亦悟。乃厚遇齐使。使张子房卒立信为齐王。封平以户牖乡。用其奇计策。卒灭楚。常以护军中尉从定燕王臧荼。\\
\hspace*{2em}汉六年。人有上书告楚王韩信反。高帝问诸将。诸将曰。﹁亟发兵阬竖子耳。﹂高帝默然。问陈平。平固辞谢。曰。﹁诸将云何。﹂上具告之。陈平曰。﹁人之上书言信反。有知之者乎。﹂曰。﹁未有。﹂曰。﹁信知之乎。﹂曰。﹁不知。﹂陈平曰。﹁陛下精兵孰与楚。﹂上曰。﹁不能过。﹂平曰。﹁陛下将用兵有能过韩信者乎。﹂上曰。﹁莫及也。﹂平曰。﹁今兵不如楚精。而将不能及。而举兵攻之。是趣之战也。窃为陛下危之。﹂上曰。﹁为之柰何。﹂平曰。﹁古者天子巡狩。会诸侯。南方有云梦。陛下弟出伪游云梦。会诸侯於陈。陈。楚之西界。信闻天子以好出游。其势必无事而郊迎谒。谒。而陛下因禽之。此特一力士之事耳。﹂高帝以为然。乃发使告诸侯会陈。﹁吾将南游云梦﹂。上因随以行。行未至陈。楚王信果郊迎道中。高帝豫具武士。见信至。即执缚之。载後车。信呼曰。﹁天下已定。我固当烹。﹂高帝顾谓信曰。﹁若毋声。而反。明矣。﹂武士反接之。遂会诸侯于陈。尽定楚地。还至雒阳。赦信以为淮阴侯。而与功臣剖符定封。\\
\hspace*{2em}於是与平剖符。世世勿绝。为户牖侯。平辞曰。﹁此非臣之功也。﹂上曰。﹁吾用先生谋计。战胜剋敌。非功而何。﹂平曰。﹁非魏无知臣安得进。﹂上曰。﹁若子可谓不背本矣。﹂乃复赏魏无知。其明年。以护军中尉从攻反者韩王信於代。卒至平城。为匈奴所围。七日不得食。高帝用陈平奇计。使单于阏氏。围以得开。高帝既出。其计礻必。世莫得闻。\\
\hspace*{2em}高帝南过曲逆。上其城。望见其屋室甚大。曰。﹁壮哉县。吾行天下。独见洛阳与是耳。﹂顾问御史曰。﹁曲逆户口几何。﹂对曰。﹁始秦时三万馀户。间者兵数起。多亡匿。今见五千户。﹂於是乃诏御史。更以陈平为曲逆侯。尽食之。除前所食户牖。\\
\hspace*{2em}其後常以护军中尉从攻陈豨及黥布。凡六出奇计。辄益邑。凡六益封。奇计或颇祕。世莫能闻也。\\
\hspace*{2em}高帝从破布军还。病创。徐行至长安。燕王卢绾反。上使樊哙以相国将兵攻之。既行。人有短恶哙者。高帝怒曰。﹁哙见吾病。乃冀我死也。﹂用陈平谋而召绛侯周勃受诏床下。曰。﹁陈平亟驰传载勃代哙将。平至军中即斩哙头。﹂二人既受诏。驰传未至军。行计之曰。﹁樊哙。帝之故人也。功多。且又乃吕后弟吕嬃之夫。有亲且贵。帝以忿怒故。欲斩之。则恐後悔。宁囚而致上。上自诛之。﹂未至军。为坛。以节召樊哙。哙受诏。即反接载槛车。传诣长安。而令绛侯勃代将。将兵定燕反县。\\
\hspace*{2em}平行闻高帝崩。平恐吕太后及吕嬃谗怒。乃驰传先去。逢使者诏平与灌婴屯於荥阳。平受诏。立复驰至宫。哭甚哀。因奏事丧前。吕太后哀之。曰。﹁君劳。出休矣。﹂平畏谗之就。因固请得宿卫中。太后乃以为郎中令。曰。﹁傅教孝惠。﹂是後吕嬃谗乃不得行。樊哙至。则赦复爵邑。\\
\hspace*{2em}孝惠帝六年。相国曹参卒。以安国侯王陵为右丞相。陈平为左丞相。\\
\hspace*{2em}王陵者。故沛人。始为县豪。高祖微时。兄事陵。陵少文。任气。好直言。及高祖起沛。入至咸阳。陵亦自聚党数千人。居南阳。不肯从沛公。及汉王之还攻项籍。陵乃以兵属汉。项羽取陵母置军中。陵使至。则东乡坐陵母。欲以招陵。陵母既私送使者。泣曰。﹁为老妾语陵。谨事汉王。汉王。长者也。无以老妾故。持二心。妾以死送使者。﹂遂伏剑而死。项王怒。烹陵母。陵卒从汉王定天下。以善雍齿。雍齿。高帝之仇。而陵本无意从高帝。以故晚封。为安国侯。\\
\hspace*{2em}安国侯既为右丞相。二岁。孝惠帝崩。高后欲立诸吕为王。问王陵。王陵曰。﹁不可。﹂问陈平。陈平曰。﹁可。﹂吕太后怒。乃详迁陵为帝太傅。实不用陵。陵怒。谢疾免。杜门竟不朝请。七年而卒。\\
\hspace*{2em}陵之免丞相。吕太后乃徙平为右丞相。以辟阳侯审食其为左丞相。左丞相不治。常给事於中。\\
\hspace*{2em}食其亦沛人。汉王之败彭城西。楚取太上皇。吕后为质。食其以舍人侍吕后。其後从破项籍为侯。幸於吕太后。及为相。居中。百官皆因决事。\\
\hspace*{2em}吕嬃常以前陈平为高帝谋执樊哙。数谗曰。﹁陈平为相非治事。日饮醇酒。戏妇女。﹂陈平闻。日益甚。吕太后闻之。私独喜。面质吕嬃於陈平曰。﹁鄙语曰﹃兒妇人口不可用﹄。顾君与我何如耳。无畏吕嬃之谗也。﹂\\
\hspace*{2em}吕太后立诸吕为王。陈平伪听之。及吕太后崩。平与太尉勃合谋。卒诛诸吕。立孝文皇帝。陈平本谋也。审食其免相。\\
\hspace*{2em}孝文帝立。以为太尉勃亲以兵诛吕氏。功多。陈平欲让勃尊位。乃谢病。孝文帝初立。怪平病。问之。平曰。﹁高祖时。勃功不如臣平。及诛诸吕。臣功亦不如勃。原以右丞相让勃。﹂於是孝文帝乃以绛侯勃为右丞相。位次第一。平徙为左丞相。位次第二。赐平金千斤。益封三千户。\\
\hspace*{2em}居顷之。孝文皇帝既益明习国家事。朝而问右丞相勃曰。﹁天下一岁决狱几何。﹂勃谢曰。﹁不知。﹂问。﹁天下一岁钱穀出入几何。﹂勃又谢不知。汗出沾背。愧不能对。於是上亦问左丞相平。平曰。﹁有主者。﹂上曰。﹁主者谓谁。﹂平曰。﹁陛下即问决狱。责廷尉。问钱穀。责治粟内史。﹂上曰。﹁苟各有主者。而君所主者何事也。﹂平谢曰。﹁主臣。陛下不知其驽下。使待罪宰相。宰相者。上佐天子理阴阳。顺四时。下育万物之宜。外镇抚四夷诸侯。内亲附百姓。使卿大夫各得任其职焉。﹂孝文帝乃称善。右丞相大惭。出而让陈平曰。﹁君独不素教我对。﹂陈平笑曰。﹁君居其位。不知其任邪。且陛下即问长安中盗贼数。君欲彊对邪。﹂於是绛侯自知其能不如平远矣。居顷之。绛侯谢病请免相。陈平专为一丞相。\\
\hspace*{2em}孝文帝二年。丞相陈平卒。谥为献侯。子共侯买代侯。二年卒。子简侯恢代侯。二十三年卒。子何代侯。二十三年。何坐略人妻。弃市。国除。\\
\hspace*{2em}始陈平曰。﹁我多阴谋。是道家之所禁。吾世即废。亦已矣。终不能复起。以吾多阴祸也。﹂然其後曾孙陈掌以卫氏亲贵戚。原得续封陈氏。然终不得。\\
\hspace*{2em}太史公曰。陈丞相平少时。本好黄帝。老子之术。方其割肉俎上之时。其意固已远矣。倾侧扰攘楚魏之间。卒归高帝。常出奇计。救纷纠之难。振国家之患。及吕后时。事多故矣。然平竟自脱。定宗庙。以荣名终。称贤相。岂不善始善终哉。非知谋孰能当此者乎。\\

% ===== 索隐述赞（司马贞《史记索隐》）· 统一样式 =====
% 题头「索隐述赞曰」五字并入反白框；正文每四句一列，定宽盒 \hfil 居列中；
% \fontsize{24}{41} 列距对齐全书界行（121+41k）；\clearpage 使述赞独占一页。
\clearpage
\shuzanhead{索隐述赞曰}\linebreak
{\shuzanbody\fontsize{24}{41}\selectfont
\hbox to \hsize{\hfil 曲逆穷巷。门多长者。宰肉先均。佐丧後罢。\hfil}\linebreak
\hbox to \hsize{\hfil 魏楚更用。腹心难假。弃印封金。刺船露裸。\hfil}\linebreak
\hbox to \hsize{\hfil 间行归汉。委质麾下。荥阳计全。平城围解。\hfil}\linebreak
\hbox to \hsize{\hfil 推陵让勃。裒多益寡。应变合权。克定宗社。\hfil}%
\par}
\end{document}
